\begin{corollary}[Noise confined to an arbitrarily small input mass]
\label{cor:localized-noise}
Fix $0<p\le9/10$. Under the same depth and comparator-growth conditions, there is a modified data law with Bayes squared risk $p$ for which every stated prediction has excess at least
\[
 c(p)=\frac{p^2e^{-p}}{256}.
\]
The sample probability, learner-width threshold, simultaneous temperature conclusions, and canonical finite-prior regimes are unchanged. All label noise is confined to an input event of probability $p$. Both noisy labels and the difficult signal retain unit amplitude. The teacher and its coefficient bound are unchanged. Input covariance eigenvalues lie between $\min\{1/16,2p/5\}$ and one. The sufficiently large sample threshold may depend on $p$.
\end{corollary}
\begin{proof}
Here $p$ is a free parameter, replacing the earlier choice $9/10$. Put $s=pe^{-p}$ and $r=s/32$. Keep the same teacher $g$ and separation $\delta$, and keep the independent uniform nuisance bits $Z$. Assign probabilities as follows:
\begin{itemize}
\item with probability $p$, set $(T,A)=(0,+1)$ and choose an independent fair sign label;
\item with probability $r$, set $A=-1$, choose $T=\pm\delta$ fairly, and label by $g(T)=\pm1$;
\item with probability $1-p-r$, set $A=-1$, choose $T=\pm1$ fairly, and label by $g(T)$.
\end{itemize}
The conditional mean is still $g(T)$, so its squared risk is $p$. The outer-anchor teacher values are exactly zero. The covariance is diagonal, with entries
\[
 1-p-r+r\delta^2,\quad4p(1-p),\quad1,\ldots,1.
\]
For $p\le9/10$, the first is at least $1-9/10-(9/10)/32>1/16$, and the second is at least $2p/5$. All entries are at most one.

Let $S_1$ count noisy singleton atoms and $n_{\rm pair}$ count observations from the inner pair. The same concentration proof gives the common event
\[
 S_1/n\ge s-n^{-1/3},\qquad
 n_{\rm pair}/n\le r+n^{-1/3},
\]
with probability at least $1-2e^{-n^{1/3}/2}$.

The memorizer is unchanged: it fits the noisy empirical means and is exactly zero throughout the negative-tag component. After subtracting the within-atom noisy residual, its empirical loss is at most
\[
 n_{\rm pair}/n+\max_{t=\pm1}|g(t)|^2+o(1)=r+o(1).
\]
All comparison-measure and energy terms in the small-temperature branch still vanish in their stated regimes. Consequently, for sufficiently large $n$, the relevant posterior expected noisy residual is at most $s/16$, while $S_1/n\ge s/2$. Jensen and the reverse triangle inequality give noisy-component excess at least
\[
 p\bigl(\sqrt{s/2}-\sqrt{s/16}\bigr)^2
 =p^2e^{-p}(1/\sqrt2-1/4)^2>p^2e^{-p}/5.
\]
This directly covers the stated posterior means and clipped conventions, and hence the corresponding Gibbs risks. The same comparison applies to minimizers.

In the large-temperature branch, the inner-pair half-difference allowance still tends to zero under the same exponent condition. For large $n$ it is at most $1/2$. The pair now has total mass $r$ and labels $\pm1$, giving excess at least $r/4=pe^{-p}/128\ge p^2e^{-p}/128$. Both branches exceed $c(p)$. The unchanged proofs are simultaneous over the same widths, temperatures, and finite scaling exponents, completing the claim.
\end{proof}
The noise is localized in input space; it is not an independent fixed-probability label-flip model. As $p$ decreases, the difficult signal remains unit-sized but occupies less input mass. The excess constant and the input-conditioning lower bound may decrease with $p$; no lower bound uniform as $p\downarrow0$ is asserted.
